\documentclass[10pt,a4paper]{amsart}
\usepackage[margin=19mm]{geometry}
\usepackage{amsmath,amssymb,mathtools}
\usepackage[scr=rsfs,cal=euler]{mathalfa}
\usepackage{xfrac}
\usepackage[unicode,hidelinks,bookmarks=false]{hyperref}
\newcommand{\ie}{i.e.\ }
\newcommand{\cf}{cf.\ }
\newcommand{\resp}{respectively\ }
\newcommand{\Subsection}{\subsection}
\providecommand{\nobreakdash}{\nobreak\hbox{-}}
\setlength{\emergencystretch}{2em}
\multlinegap0pt
\usepackage{amssymb}
\usepackage{mathtools}
\usepackage[shortlabels]{enumitem}
\usepackage{xcolor}
\newtheorem{theorem}{Theorem}
\newcommand{\C}{\mathbb C}
\newcommand{\M}{\mathbb M}
\newcommand{\TZ}{C^{\mathrm{TZ}}}
\newcommand{\Tr}{\operatorname{Tr}}
\title{A Counterexample to the Tang Zhang Schatten Norm Conjecture and Sharp Positive Results}
\author{Zijian Zeng, Houde Liu, Kurunathan Ratnavelu}
\date{Source: arXiv:2608.15558v1 (CC BY 4.0)}
\begin{document}
\maketitle

\noindent\emph{Source and licence.} A Counterexample to the Tang Zhang Schatten Norm Conjecture and Sharp Positive Results. Version arXiv:2608.15558v1 (CC BY 4.0), distributed by arXiv under the Creative Commons Attribution 4.0 International licence, http://creativecommons.org/licenses/by/4.0/, as recorded in the arXiv licence metadata for this exact version and shown on the abstract page https://arxiv.org/abs/2608.15558v1. Full licence notice: \texttt{licenses/arxiv-2608.15558-CC-BY-4.0.txt}.\par\smallskip

\noindent\textit{Editorial scope.} Counted unit: the article's theorem thm:p4 (source0.tex lines 189--202). The article's complete original proof of it is delivered with it (source0.tex lines 506--654, one \texttt{proof} environment of 149 lines, which the article places away from the statement). No mathematics is written, repaired or re-derived: the statement and the proof are the article's own lines, in the article's own notation, and no retained line is substituted or re-flowed.  No further source-local context is needed: every cross-reference of every retained line resolves inside the two delivered spans.  Nothing was omitted from any retained span, so the package carries no omission record.  No retained line contains a citation, so the wrapper carries no bibliography and no bibliography entry.  No retained line contains a figure environment, an image-inclusion command or drawing code, so no image is delivered and the package declares no asset dependency.  Editorial formatting only, in addition: an AMS \texttt{amsart} preamble that restates the article's own declarations and macros used by the retained lines, namely the environments \texttt{theorem} on separate counters, together with the standard packages those lines need.  The retained environments print in this document as Theorem 1 in order of appearance, whereas the article prints the same items with the numbering of its own full document.  The complete native source of the article as distributed in the arXiv e-print for this exact version is bundled unchanged as \texttt{sources/207758/source0.tex}.
\begin{theorem}[The full case \(m=2,\ p=4\)]\label{thm:p4}
Let \(A,B\in\M_n(\C)\), and let \(x>1\) be the solution of
\[
 x^4-2x-1=0.
\]
Then
\begin{equation}\label{eq:p4-bound}
 \|A+B\|_4^4
 \le
 \frac{x^2(x+1)}2\,\||A|+|B|\|_4^4.
\end{equation}
The constant is sharp in the dimension-free problem, is attained for every
\(n\ge2\), and equals \((\TZ_{4,2})^4\).
\end{theorem}

\begin{proof}[Proof of Theorem~\ref{thm:p4}]
Put \(H=|A|\) and \(K=|B|\).  Extend the partial isometries in the polar
decompositions to unitaries, and write \(A=UH,\ B=VK\).
With \(W=U^*V\), unitary invariance reduces the numerator to
\[
 S=H+WK.
\]
Set
\[
 a=\Tr H^4,\quad b=\Tr K^4,\quad d=\Tr H^2K^2.
\]
The Hilbert--Schmidt triangle inequality applied to
\[
 SS^*=H^2+HKW^*+WKH+WK^2W^*
\]
gives
\begin{equation}\label{eq:p4-numerator}
 \|S\|_4^4
 \le(\sqrt a+\sqrt b+2\sqrt d)^2.
\end{equation}

Introduce
\[
 e=\Tr H^3K,\qquad f=\Tr HK^3,\qquad g=\Tr HKHK.
\]
Cyclically collecting all words in the noncommutative expansion gives
\begin{equation}\label{eq:p4-expansion}
 \Tr(H+K)^4=a+b+4(e+f)+4d+2g.
\end{equation}
Two lower bounds are needed.  First,
\begin{equation}\label{eq:p4-positive}
 e+f-2d
 =
 \Tr\big((H-K)H(H-K)K\big)
 =
 \|H^{1/2}(H-K)K^{1/2}\|_2^2
 \ge0.
\end{equation}
Second, the three-factor Schatten H\"older inequality gives
\begin{align*}
 \sqrt d=\|HK\|_2
 &=
 \|H^{1/2}(H^{1/2}K^{1/2})K^{1/2}\|_2\\
 &\le
 \|H^{1/2}\|_8
 \|H^{1/2}K^{1/2}\|_4
 \|K^{1/2}\|_8
 =
 a^{1/8}g^{1/4}b^{1/8}.
\end{align*}
Thus, when \(ab>0\),
\begin{equation}\label{eq:g-lower}
 g\ge\frac{d^2}{\sqrt{ab}}.
\end{equation}
If \(ab=0\), one of \(H,K\) vanishes and the theorem is immediate.
Combining \eqref{eq:p4-expansion}--\eqref{eq:g-lower}, we obtain
\begin{equation}\label{eq:p4-denominator}
 \Tr(H+K)^4
 \ge a+b+12d+\frac{2d^2}{\sqrt{ab}}.
\end{equation}

Let \(A_0=\sqrt a,\ B_0=\sqrt b\), and define
\[
 z=\frac{A_0+B_0}{\sqrt{A_0B_0}}\ge2,\qquad
 s=\sqrt{\frac{d}{A_0B_0}}\in[0,1].
\]
The upper bound on \(s\) is the Hilbert--Schmidt Cauchy--Schwarz
inequality \(d\le\sqrt{ab}\).  From
\eqref{eq:p4-numerator} and \eqref{eq:p4-denominator},
\begin{equation}\label{eq:Phi}
 \frac{\|A+B\|_4^4}{\||A|+|B|\|_4^4}
 \le
 \Phi(z,s)
 :=
 \frac{(z+2s)^2}{z^2-2+12s^2+2s^4}.
\end{equation}

If \(s\ge1/2\), then
\[
 2(z^2-2+12s^2+2s^4)-(z+2s)^2
 =
 (z-2s)^2+16s^2+4s^4-4\ge0,
\]
so \(\Phi(z,s)\le2\).

Suppose \(0\le s\le1/2\).  The sign of
\(\partial\Phi/\partial z\) is the sign of
\[
 s^4+6s^2-zs-1.
\]
Since \(z\ge2\),
\[
 s^4+6s^2-zs-1
 \le s^4+6s^2-2s-1
 \le2s-1\le0;
\]
the middle inequality follows from
\(s^3+6s-4\le1/8+3-4<0\).
Consequently,
\begin{equation}\label{eq:f-s}
 \Phi(z,s)\le\Phi(2,s)
 =
 f(s):=\frac{2(1+s)^2}{1+6s^2+s^4}.
\end{equation}
The sign of \(f'(s)\) is the sign of
\[
 1-6s-2s^3-s^4.
\]
Hence \(f\) has a unique maximizer \(s_0\in(0,1/2)\), characterized by
\begin{equation}\label{eq:s0}
 s_0^4+2s_0^3+6s_0-1=0.
\end{equation}
The first branch cannot dominate, because \(f(1/6)>2\).

Set \(x=(1+s_0)/(1-s_0)\).  A direct calculation gives
\[
 (x^4-2x-1)(1-s_0)^4
 =
 2(s_0^4+2s_0^3+6s_0-1)=0,
\]
and
\begin{equation}\label{eq:f-constant}
 f(s_0)
 =
 \frac{x^2(x+1)^2}{x^4+1}
 =
 \frac{x^2(x+1)}2,
\end{equation}
where the last equality uses \(x^4+1=2(x+1)\).
Equations~\eqref{eq:Phi}--\eqref{eq:f-constant} prove
\eqref{eq:p4-bound}.

For sharpness, choose unit vectors \(r_1,r_2\) with
\(\langle r_1,r_2\rangle=s_0\), choose a unit vector \(u\), and set
\[
 A=ur_1^*,\qquad B=ur_2^*.
\]
Then
\[
 \|A+B\|_4^4=4(1+s_0)^2
\]
and
\[
 \||A|+|B|\|_4^4
 =(1+s_0)^4+(1-s_0)^4
 =2(1+6s_0^2+s_0^4).
\]
Their ratio is \(f(s_0)\), proving sharpness.
\end{proof}

\end{document}
